\section{Operacioni pojačavač u zatvorenoj sprezi}

{Problem kod izvedbe operacionog pojačavača u otvorenoj sprezi je u tome što se izlazni napon $V_{out}$ veoma lako može dovesti u stanje saturacije. Zbog veoma velikog pojačanja $A_{o}$, čak i izuzetno male vrednosti diferencijalnog napona na ulazu mogu dovesti do maksimalne vrednosti izlaznog napona. Zbog toga je operacioni pojačavač u otvorenoj sprezi praktično neupotrebljiv za potrebe linearnog pojačanja signala. Na primer, ako bi diferencijalni napon na ulazu iznosio $1,\mu\mathrm{V}$, na izlazu bi se dobila vrednost od $100,\mathrm{V}$. Međutim, isti izlazni napon bi se dobio i za ulazne napone od $10,\mu\mathrm{V}$, $100,\mu\mathrm{V}$, $1,\mathrm{mV}$ itd. Na taj način praktično nije moguće razlikovati različite vrednosti ulaznog signala, jer sve one rezultuju istom vrednošću izlaznog signala, odnosno ne dobijamo linearno mapiranje ulaznog signala na izlazni signal.}

{Zbog toga se operacioni pojačavači u praksi koriste u zatvorenoj sprezi, pri čemu se deo izlaznog signala vraća na ulaz operacionog pojačavača. Na taj način se postiže rad operacionog pojačavača u linearnom režimu, odnosno omogućava se da se izlazni signal $V_{out}$ menja u skladu sa ulaznim signalom $V_{in}$, pri čemu se izbegava saturacija. Ovime se dobija linearno mapiranje ulaznog signala na izlazni.}

{Postoje dve izvedbe operacionog pojačavača u zatvorenoj sprezi koje se najčešće koriste, a to su: neinvertujući i invertujući operacioni pojačavač.}

\subsection{Neinvertujući operacioni pojačavač}

\begin{figure}[ht]
    \centering
    \begin{circuitikz}[american]
        \begin{scope}
            \ctikzset{tripoles/op amp/width=2.5,   % default 2.5
                    tripoles/op amp/height=2}    % default 2
            \node[op amp] (OA) at (0,0) {};
        \end{scope}
        \draw
            (OA.+) -- ++(-2,0) to [V, l_=$V_{in}$] ++(0,-1) node[ground] {}
            (OA.-) to [R, l_=$R_1$] ++(-4, 0) node[ground] {}
            (OA.out) -- ++(1,0) node[right]{$V_{out}$};

        \draw
            (OA.up) -- ++(0,1) node[above]{$+V_{cc}$};

        \draw
            (OA.down) -- ++(0,-1) node[below]{$-V_{cc}$};
        
        \coordinate (topV_minus) at ($(OA.-) + (0, 2.5)$);
        \coordinate (startarrow) at ($(OA.-) + (-1.25, 0)$);

        \draw[->]
            (startarrow) -- ++(1,0) node[midway, above] {$I_1$};

        \draw[->]
            (topV_minus) -- ++(1,0) node[midway, above] {$I_2$};

        \node[circ] at (OA.-) {};
        \node[below] at (OA.-) {$V^-$};
        \node[circ] at (OA.+) {};
        \node[below] at (OA.+) {$V^+$};
        \node[circ] at (OA.out) {};

        \draw
            (OA.-) -- (topV_minus) to [R, l=$R_2$] ++(3.5,0) -- (OA.out);

        
    \end{circuitikz}
    \caption{Neinvertujući operacioni pojačavač}
    \label{fig:opamp_noninverting}
\end{figure}

{U konfiguraciji neinvertujećeg operacionog pojačavača ulazni signal $V_{in}$ povezan je na neinvertujući ulaz $V^+$ operacionog pojačavača, a invertujući ulaz se preko otpornika $R_1$ dovodi na masu. Pod pretpostavkom da se radi o idealnom operacionom pojačavaču važi da je}

\begin{equation}
    V^- = V^+ = V_{in}.
\end{equation}

{Budući da je ulazna impedansa idealnog operacionog pojačavača beskonačna, kroz njegove ulazne priključke ne teče struja. Zbog toga su struje kroz otpornike $R_1$ i $R_2$ jednake, pa važi}

\begin{equation}
    I_1 = \frac{V^- - 0}{R_1} \quad i \quad I_2 = \frac{V_{out} - V^-}{R_2},
\end{equation}

{odnosno}

\begin{equation}
    I = \frac{V^- - 0}{R_1} = \frac{V_{out} - V^-}{R_2}.
    \label{eq:opamp_noninv}
\end{equation}

{Nakon izjednačavanja jednačina u izrazu \ref{eq:opamp_noninv} dobija se izraz za zavisnost izlaznog napona operacionog pojačavača od ulaznog napona}

\begin{equation}
    V_{out} = (1 + \frac{R_2}{R_1})V_{in},
\end{equation}

{gde je izraz $1 + \frac{R_2}{R_1}$ zapravo pojačanje operacionog pojačavača, odnosno}

\begin{equation}
    A_o = 1 + \frac{R_2}{R_1}.
\end{equation}

{Za razliku od otvorene sprege, gde je pojačanje određeno karakteristikama samog operacionog pojačavača, u konfiguraciji sa negativnom povratnom spregom pojačanje je određeno vrednostima otpornika $R_1$ i $R_2$.}

\subsection{Invertujući operacioni pojačavač}

\begin{figure}[ht]
    \centering
    \begin{circuitikz}[american]
        \begin{scope}
            \ctikzset{tripoles/op amp/width=2.5,   % default 2.5
                    tripoles/op amp/height=2}    % default 2
            \node[op amp] (OA) at (0,0) {};
        \end{scope}
        \draw
            (OA.+) -- ++(-2,0) node[ground] {}
            (OA.-) to [R, l_=$R_1$] ++(-4, 0) to [V, l_=$V_{in}$] ++(0,-1.4) node[ground] {}
            (OA.out) -- ++(1,0) node[right]{$V_{out}$};

        \draw
            (OA.up) -- ++(0,1) node[above]{$+V_{cc}$};

        \draw
            (OA.down) -- ++(0,-1) node[below]{$-V_{cc}$};
        
        \coordinate (topV_minus) at ($(OA.-) + (0, 2.5)$);
        \coordinate (startarrow) at ($(OA.-) + (-1.25, 0)$);

        \draw[->]
            (startarrow) -- ++(1,0) node[midway, above] {$I_1$};

        \draw[->]
            (topV_minus) -- ++(1,0) node[midway, above] {$I_2$};

        \node[circ] at (OA.-) {};
        \node[below] at (OA.-) {$V^-$};
        \node[circ] at (OA.+) {};
        \node[below] at (OA.+) {$V^+$};
        \node[circ] at (OA.out) {};

        \draw
            (OA.-) -- (topV_minus) to [R, l=$R_2$] ++(3.5,0) -- (OA.out);

        
    \end{circuitikz}
    \caption{Invertujući operacioni pojačavač}
    \label{fig:opamp_inverting}
\end{figure}

{U ovoj konfiguraciji operacionog pojačavača se neinvertujući ulaz $V^+$ povezuje na masu, a na invertujući ulaz je povezano ulazno napajanje $V_{in}$ preko otpornika $R_1$. Pod pretpostavkom da se radi o idealnom operacionom pojačavaču može da se zaključi da su tačke $V^-$ i $V^+$ na istom potencijalu, odnosno da je $V^-$ povezan na virtuelnu masu preko $V^+$ i važi}

\begin{equation}
    V^- = V^+ = 0
\end{equation}

{Kao i kod izvedbe neinvertujućeg operacionog pojačavača i ovde je izraz za struju u kolu}

\begin{equation}
    I_1 = \frac{V_{in} - 0}{R_1} \quad i \quad I_2 = \frac{0 - V_{out}}{R_2},
\end{equation}

{odnosno}

\begin{equation}
    I = \frac{V_{in} - 0}{R_1} = \frac{0 - V_{out}}{R_2}.
\end{equation}

{I odavde može da se zaključi da je izraz koji određuje zavisnost izlaza od ulaza jednak}

\begin{equation}
    V_{out} = -\frac{R_2}{R_1} \cdot V_{in}.
\end{equation}

{Treba primetiti da se u ovoj izvedbi operacionog pojačavača ulazni signal invertuje i u zavisnosti od vrednosti otpornika $R_1$ i $R_2$, signal takođe može biti pojačan, oslabljen ili samo biti invertovan.}